\section{Introduction}

\subsection{Elevator Pitch}

% One or two sentences. What is the game, what does the player do,
% and what is the single most distinctive thing about it?
% Example shape: "<TITLE> is a <GENRE> in which the player
% <CORE VERB> in order to <GOAL>, distinguished by <HOOK>."

\textit{[One- or two-sentence pitch goes here.]}

\subsection{Concept}

% A short paragraph (4--8 lines) expanding the pitch. Cover:
%   - Setting / tone / mood.
%   - The player's role and primary fantasy.
%   - The central tension or conflict.
%   - Why this game exists --- what question or feeling it is chasing.

\textit{[Concept paragraph goes here.]}

\subsection{Design Pillars}

% 3--5 short, opinionated statements that every design decision is
% measured against. Pillars should be specific enough to reject
% ideas, not just describe taste.

\begin{itemize}
    \item \textbf{Pillar 1.} \textit{[Short statement.]}
    \item \textbf{Pillar 2.} \textit{[Short statement.]}
    \item \textbf{Pillar 3.} \textit{[Short statement.]}
\end{itemize}

\subsection{Genre and Inspirations}

% Name the genre(s) plainly. List 3--5 reference works (games, films,
% books) and, for each, a single sentence on what is being borrowed
% or deliberately diverged from.

\begin{itemize}
    \item \textit{[Reference]} --- \textit{[what is borrowed / what is rejected].}
    \item \textit{[Reference]} --- \textit{[\ldots].}
\end{itemize}

\subsection{Core Loop}

% Describe the moment-to-moment loop in 3--5 steps. This is what
% the player does over and over in a single play session.

\begin{enumerate}
    \item \textit{[Step 1 --- the player \ldots]}
    \item \textit{[Step 2 --- which leads to \ldots]}
    \item \textit{[Step 3 --- and finally \ldots]}
\end{enumerate}

\subsection{Scope of this Document}

This is a living document. Sections evolve as the project does.

\begin{itemize}
    \item \textbf{Mechanics} --- design rationale and implementation
          notes for the game's core systems.
    \item \textbf{Devlog} --- chronological entries describing
          progress, problems encountered, and decisions made.
\end{itemize}

\subsection{Project Conventions}

% Replace with whatever is true for this project. Keep it short ---
% this section is a reference card, not documentation.

\begin{itemize}
    \item \textbf{Engine:} \textit{[engine and version]}.
    \item \textbf{Language:} \textit{[language / runtime]}.
    \item \textbf{Source layout:} \textit{[where gameplay code lives,
          where scenes/prefabs live]}.
    \item \textbf{Naming:} \textit{[namespace / folder conventions]}.
\end{itemize}

\newpage
